\chapter{Types}
\label{ch:types}
\index{types}

Every expression has a type, and a program is only run once all its types
check. Types are inferred: annotations are optional, and when present they
are checked against the inferred type.

\section{The types}

\begin{center}
\small
\begin{tabularx}{\textwidth}{@{}l X@{}}
\toprule
\textbf{Type} & \textbf{Values} \\
\midrule
\tc{Int} & machine integers (\S\ref{sec:int-width}) \\
\tc{Float} & double-precision floating point numbers \\
\tc{Bool} & \tc{true}, \tc{false} \\
\tc{Char} & single characters \\
\tc{String} & immutable strings \\
\tc{Unit} & the single value \tc{()} \\
\tc{a -> b} & functions from \tc{a} to \tc{b}; \tc{->} associates to the right \\
\tc{List<a>} (or \tc{[a]}) & lists whose elements have type \tc{a} \\
\tc{(a, b, c)} & tuples \\
\tc{{ x: Int, y: Int }} & records \\
\tc{Option<a>}, \tc{Ordering}, \dots & sum types, declared with \tc{type} \\
\bottomrule
\end{tabularx}
\end{center}

\noindent Lower-case letters in a type --- \fc|a|, \fc|b| --- are \emph{type
variables}: \fc|a -> a| is the type of a function that returns a value of the
same type as its argument, whatever that type is.

\section{Type declarations}
\label{sec:adts}
\index{sum types}\index{type@\texttt{type}}

\fc@type Name<a, b> = | C1 of T1 | C2 | ...@ declares a \emph{sum type}: a
value is one of the constructors, with a payload of the given type or none.
A constructor with a payload is a function: \fc|Some : a -> Option<a>|; a
constructor with several fields takes a tuple, \fc|Node (l, v, r)|.
A type that refers to itself is declared with \fc|type rec|, and several
types that refer to each other with \fc|type rec ... and ...|.
\fc|type Name = T|, where \fc|T| is not a list of constructors, declares an
\emph{alias}.

\forgeexample{types-adts}

\forgeexample{types-mutual}

\section{Annotations}
\index{annotations}
A pattern, a parameter or a result can be annotated: \fc|(x : Int)|,
\fc|let f (x : Int) : String = ...|, \fc|let origin : Pair<Int> = (0, 0)|.
Type variables in an annotation stand for any type that makes the program
check, so \fc|(f : a -> a)| merely says that \fc|f| returns what it takes.

\forgeexample{types-annotations}

\section{Records}
\label{sec:records}
\index{records!row polymorphism}
Record types are \emph{structural}: two records with the same fields have the
same type, whether or not the type has a name. A function that uses some
fields of a record accepts every record that has them --- its inferred type
has a \emph{row variable} standing for ``the other fields''. A record
literal, on the other hand, has exactly its fields, so a field that does not
exist is a type error.

\forgeexample{types-rows}

\section{Type inference and polymorphism}
\label{sec:inference}
\index{type inference}\index{polymorphism}

Forge infers types with the Hindley--Milner algorithm, extended with records
and traits. A name bound by \fc|let| (at top level or locally) gets the most
general type that fits its definition, and each use of the name can
instantiate its type variables differently. Lambda-bound names (function
parameters) are not generalised: a parameter has one type inside its
function.

\forgeexample{types-poly}

Section~\ref{sec:static-semantics} gives the typing rules, and
Section~\ref{sec:impl-inference} describes the algorithm.
